% !TeX program = pdflatex
\documentclass[11pt,a4paper]{article}
\usepackage[T1]{fontenc}
\usepackage[utf8]{inputenc}
\usepackage{lmodern,amsmath,amssymb,amsthm,mathtools,mathrsfs,microtype}
\usepackage[margin=23mm]{geometry}
\usepackage[hidelinks]{hyperref}
\usepackage{xurl}
\setlength{\emergencystretch}{3em}
\allowdisplaybreaks[2]
\newtheorem{theorem}{Theorem}[section]
\newtheorem{proposition}[theorem]{Proposition}
\newtheorem{lemma}[theorem]{Lemma}
\newtheorem{corollary}[theorem]{Corollary}
\theoremstyle{remark}\newtheorem{remark}[theorem]{Remark}
\newcommand{\Ric}{\operatorname{Ric}}
\newcommand{\Rm}{\operatorname{Rm}}
\newcommand{\tr}{\operatorname{tr}}
\newcommand{\diver}{\operatorname{div}}
\newcommand{\Met}{\operatorname{Met}}
\newcommand{\Vol}{\operatorname{Vol}}
\newcommand{\CP}{\mathbb{CP}}
\newcommand{\ip}[2]{\langle #1,#2\rangle}
\numberwithin{equation}{section}
\title{A genuine descent direction on the harmonic-alignment branch\\[3pt]
\large Trace-free Ricci criticality and almost-K\"ahler equality geometry}
\author{Qiushi research continuation}
\date{8 October 2026}
\begin{document}
\maketitle
\begin{abstract}
We study the unrestricted functional $P(g)=\int|\Ric_g^\circ|^2$ on closed
four-manifolds. For a constant-scalar almost-K\"ahler metric with
$J$-invariant Ricci tensor, we derive the exact first variation
$DP_g[\Ric_g^\circ]=-sP(g)/2$. Thus every non-Einstein member with
$s\ne0$ has an explicit volume-preserving descending path, without a
K\"ahler, symmetry, period, or spectral assumption. Applied to the
harmonic-flux equality theorem in the supplied October 8 manuscript,
this excludes its entire non-Einstein equality branch from being a
smooth global $P$-minimizer. A related identity retains the amplitude
correlation in the more general harmonic-product case. These results
provide an additional obstruction to proposed minimizers; they do not
establish smooth attainment on the Godeaux blowup or solve the general
non-Bach-flat minimizer problem.
\end{abstract}

\section{Functional and conventions}
Let $(M,g)$ be a smooth closed connected oriented four-manifold and put
\[
 Z=\Ric-\frac{s}{4}g,\qquad P(g)=\int_M|Z|^2\,dV_g.
\]
Two-forms have the form norm, so $|e^{12}+e^{34}|^2=2$; they are
identified with skew endomorphisms when composed. Weyl curvature is an
endomorphism of $\Lambda^2$. Write
\[
 \mathscr R(H,H)=R_{ikjl}H^{ij}H^{kl},\qquad
 R_{ikjl}g^{kl}=\Ric_{ij},
\]
with positive sectional curvature convention. The rough Laplacian
$\nabla^*\nabla$ and scalar Laplacian $\Delta_0=d^*d$ are nonnegative.
A constant-scalar metric has $\diver Z=0$ by contracted Bianchi.

The classical gradients of $\int|\Ric|^2$ and $\int s^2$ in
\cite[(2.5), (2.9)]{GV} give, at constant scalar curvature,
\begin{equation}\label{eq:gradient}
 \mathcal G_P=\nabla^*\nabla Z-2\mathscr R(Z)+\tfrac12|Z|^2g.
\end{equation}
Here $DP_g[H]=\int\ip{\mathcal G_P}{H}\,dV_g$. Consequently
\begin{equation}\label{eq:variation}
 DP_g[Z]=\int_M\bigl(|\nabla Z|^2-2\mathscr R(Z,Z)\bigr)\,dV_g.
\end{equation}
The direction $Z$ is fixed at the base metric for this derivative.
No differentiation of a choice of a field along a metric family is used.
The conformal trace equation for a closed $P$-critical metric is
$\Delta_0s=0$, so every such metric has constant scalar curvature.
Scale invariance makes criticality on unit-volume metrics equivalent
to unrestricted criticality.

\section{The opposite-duality product identity}

\begin{lemma}\label{lem:algebra}
Let $\phi\in\Lambda^+$, $\psi\in\Lambda^-$, and $H=\phi\circ\psi$.
Then $H$ is symmetric and trace-free and
\begin{align}
 |H|^2&=|\phi|^2|\psi|^2,\label{eq:isometry}\\
 \mathscr R(H,H)&=
 |\psi|^2\ip{W^+\phi}{\phi}
 +|\phi|^2\ip{W^-\psi}{\psi}
 -\frac{s}{12}|\phi|^2|\psi|^2.\label{eq:curvature}
\end{align}
For smooth forms put $f=|\phi|^2$ and $h=|\psi|^2$. Then
\begin{equation}\label{eq:derivatives}
 |\nabla H|^2=h|\nabla\phi|^2+f|\nabla\psi|^2
                  +\tfrac12\ip{df}{dh}.
\end{equation}
\end{lemma}
\begin{proof}
Opposite-duality skew endomorphisms commute. Their composition realizes
an isometry $\Lambda^+\otimes\Lambda^-\simeq S^2_0$ with these norms.
This proves \eqref{eq:isometry} and, after differentiating the product,
\eqref{eq:derivatives}. Explicitly the cross term is
$2\sum_i\ip{\nabla_i\phi}{\phi}\ip{\psi}{\nabla_i\psi}$.

At a point with both forms nonzero, choose an oriented orthonormal frame
such that $\phi=a(e^{12}+e^{34})$ and
$\psi=b(e^{12}-e^{34})$. Then
$H=ab\operatorname{diag}(-1,-1,1,1)$ and
\[
 \mathscr R(H,H)
 =2a^2b^2(K_{12}+K_{34}-K_{13}-K_{14}-K_{23}-K_{24}).
\]
Using $s=2\sum_{i<j}K_{ij}$ and the two diagonal curvature blocks
$W^\pm+sI/12$ gives \eqref{eq:curvature}. The trace-free Ricci block
drops out of this contraction; it has not been assumed zero.
The identities extend over zero sets by continuity.
\end{proof}

\begin{proposition}[Harmonic-product first variation]\label{prop:product}
Suppose $s$ is constant and there are smooth harmonic forms
$\phi\in\Omega^+(M)$ and $\psi\in\Omega^-(M)$ such that
$Z=\phi\circ\psi$. Then
\begin{equation}\label{eq:main-general}
 \boxed{DP_g[Z]=-\frac{s}{2}P(g)
           -\frac12\int_M\ip{d|\phi|^2}{d|\psi|^2}\,dV_g.}
\end{equation}
In particular, if either harmonic form has constant length, then
$DP_g[Z]=-sP(g)/2$.
\end{proposition}
\begin{proof}
The harmonic two-form Weitzenb\"ock equations are
\[
 \nabla^*\nabla\phi=2W^+\phi-\frac{s}{3}\phi,
 \qquad
 \nabla^*\nabla\psi=2W^-\psi-\frac{s}{3}\psi.
\]
Thus
\[
 |\nabla\phi|^2=2\ip{W^+\phi}{\phi}-\frac{s}{3}f-\tfrac12\Delta_0f,
 \quad
 |\nabla\psi|^2=2\ip{W^-\psi}{\psi}-\frac{s}{3}h-\tfrac12\Delta_0h.
\]
Substitute these in \eqref{eq:derivatives}, subtract twice
\eqref{eq:curvature}, and integrate. The scalar term is $-sfh/2$.
The amplitude terms are
\[
 -\tfrac12\int(h\Delta_0f+f\Delta_0h)
       +\tfrac12\int\ip{df}{dh}
 =-\tfrac12\int\ip{df}{dh}.
\]
Equation \eqref{eq:variation} now proves the assertion.
All fields are smooth on the original closed manifold, so there are
no puncture or end terms in this integration.
\end{proof}

For a nonzero negative-scalar $P$-critical metric in this general
harmonic-product class, \eqref{eq:main-general} requires
\begin{equation}\label{eq:amplitude}
 \int_M\ip{d|\phi|^2}{d|\psi|^2}\,dV_g=|s|P(g)>0.
\end{equation}
This is a necessary condition, not an exclusion of every
Einstein--Maxwell geometry. When neither amplitude is constant, its
correlation cannot be omitted.

\section{An unrestricted descent for the almost-K\"ahler branch}

\begin{theorem}\label{thm:AK}
Let $(M^4,g,J,\Omega)$ be closed almost-K\"ahler, with constant
scalar curvature $s$ and $J$-invariant Ricci tensor. Then
\begin{equation}\label{eq:AK-variation}
 \boxed{DP_g[Z]=-\frac{s}{2}P(g).}
\end{equation}
If $s\ne0$ and $g$ is not Einstein, it is not even a stationary point
of the unrestricted $P$ functional. An explicit descending family is
\begin{equation}\label{eq:path}
 g_t(X,Y)=g\bigl(\exp(t\operatorname{sgn}(s)Z^\sharp)X,Y\bigr),
 \qquad t>0\text{ small}.
\end{equation}
It preserves the volume form pointwise and satisfies
\begin{equation}\label{eq:descent}
 P(g_t)=P(g)-\frac{|s|}{2}P(g)t+O_g(t^2).
\end{equation}
Moreover
\begin{equation}\label{eq:gradient-lower}
 \|\mathcal G_P\|_{L^2}\ge\frac{|s|}{2}\sqrt{P(g)}.
\end{equation}
\end{theorem}
\begin{proof}
Use the convention $\Omega^\sharp=J$; changing its sign has no effect
on the argument. The symmetric endomorphism $Z^\sharp$ commutes with
$J$ and is trace-free. Therefore $\psi=-JZ^\sharp$ is an anti-self-dual
two-form and $Z=\Omega\circ\psi$.

The almost-K\"ahler form is self-dual, closed, and of constant norm
$\sqrt2$, hence harmonic. For forms of opposite duality the product
rule, when the first form is closed, gives
\[
 \diver(\Omega\circ\psi)=\Omega(d^*\psi).
\]
Indeed $d^*\Omega=0$; the remaining differentiated anti-self-dual
endomorphisms commute with $\Omega$, yielding the displayed identity.
Contracted Bianchi gives $\diver Z=0$. Since $J$ is invertible,
$d^*\psi=0$; anti-self-duality gives $d\psi=0$. Thus $\psi$ is
harmonic globally, including at its zeros. This also recovers the
usual harmonic primitive-Ricci-form fact without assuming that $J$
is integrable.

Apply Proposition~\ref{prop:product} with $\phi=\Omega$, whose norm
is constant. The derivative of \eqref{eq:path} at zero is
$\operatorname{sgn}(s)Z$. The endomorphism exponential is positive,
and its determinant is one because $\tr Z=0$. This proves
\eqref{eq:descent} and volume preservation. The family need not remain
in the almost-K\"ahler subspace: Q3 asks about all metrics.
Finally, Cauchy--Schwarz applied to
$|DP_g[Z]|=(|s|/2)P(g)$ proves \eqref{eq:gradient-lower}.
\end{proof}

\begin{corollary}\label{cor:critical}
A smooth closed $P$-critical metric which admits an almost-K\"ahler
structure with invariant Ricci tensor is Einstein if $s\ne0$.
If $s=0$, it is Bach-flat. In particular no metric in this class
answers the non-Bach-flat global-minimizer problem.
\end{corollary}
\begin{proof}
Criticality first makes $s$ constant. For $s\ne0$,
\eqref{eq:AK-variation} forces $P=0$.
For $s=0$, write
\[
 P=\left(\int|\Ric|^2-\frac13\int s^2\right)
                    +\frac1{12}\int s^2.
\]
The last functional has zero gradient at $s=0$. The first is,
up to a topological term and a nonzero factor, the Weyl functional.
Thus criticality is Bach-flatness \cite[Section 1.4]{GV}.
\end{proof}

\section{Application to the supplied harmonic-flux equality theorem}
For a nonzero harmonic anti-self-dual form $h$, set
\[
 H=\|h\|_2^2,\quad \mathfrak F_h=\sqrt2\int_M|B^*h|\,dV_g,
 \quad L_h=\frac{P/2-\mathfrak F_h^2/H}{4\pi^2},
 \qquad |Z|^2=4|B|^2.
\]
The supplied manuscript \cite[Section 3.2]{Source} proves that at
constant scalar curvature, $L_h=0$ implies either Einstein geometry,
or a global almost-K\"ahler structure with invariant Ricci tensor and
primitive Ricci form a constant multiple of $h$. Its bounded harmonic
form extends across the zeros of $h$ by the stated removability argument.

Combining that result with Theorem~\ref{thm:AK}, rather than solving
the negative Blair--Ianus integrability problem, gives
\begin{equation}\label{eq:equality-exclusion}
 \boxed{s\ne0,\ P(g)>0,\ L_h(g)=0
     \quad\Longrightarrow\quad
     g\text{ is not an unrestricted local }P\text{-minimizer}.}
\end{equation}
In the normalization $s=-4$, the actual path $g_t=g(e^{-tZ^\sharp}\cdot,\cdot)$
has $DP_g[-Z]=-2P(g)<0$. Replacing $g_t$ by its negative Yamabe
representative, and rescaling if desired, cannot increase $P$; hence
the decrease is not an artifact of leaving the constant-scalar locus.

If a non-Einstein, negative-scalar smooth global minimizer $g_*$ exists,
then for every nonzero harmonic ASD form $h$ one must have $L_h(g_*)>0$.
In fact, provided $b^->0$,
\begin{equation}\label{eq:fixed-gap}
 \min_{\substack{h\in\mathcal H^-(g_*)\\\|h\|_2=1}}L_h(g_*)>0.
\end{equation}
The sphere is compact in a fixed finite-dimensional vector space, and
$h\mapsto\int|B^*h|$ is continuous. A zero minimum would contradict
\eqref{eq:equality-exclusion}. This positive number is allowed to
depend on $g_*$; it is not uniform over degenerating geometries.

The latest product calculation in \cite[Section 4]{Source} concerns
$P/(8\pi^2)-|a_g^-|^2$ with a moving harmonic period. The present
calculation concerns $P$ itself under unrestricted metric variations.
These stationarity conditions are different. In particular the
new conclusion does not contradict stationarity of the corrected
functional at a non-Einstein cscK metric.

\section{The precise Q3 conclusion}
Keep the Godeaux blowup $M=X\#\overline{\CP}^{\,2}$ used in the
previous continuation, with the inherited estimates
\[
  \frac{16\pi^2}{3}\le P_*(M)<8\pi^2,
\]
and negative Yamabe constant for every conformal class. If $P_*(M)$
is attained smoothly, the attaining metric is non-Einstein and
non-Bach-flat. This reduction, the strict upper construction, and the
previous negative-lattice exclusion are project inputs \cite{Prior}.

The present result adds a definitive restriction on such an attaining
metric: it cannot belong to the entire harmonic-flux equality branch,
even if its almost-complex structure is nonintegrable. There is no need
to decide integrability of that branch before excluding it from Q3.
The proof supplies a signed, admissible metric variation, not only a
nonnegative decomposition.

It does not exclude arbitrary ALE or partially collapsed minimizing
limits. It also does not imply that a minimizing sequence has
$L_h\to0$. Smooth attainment on the fixed manifold, or another complete
example for Q3, remains to be proved. The four public questions are not
claimed solved by this note.

\begin{thebibliography}{9}
\bibitem{GV}
M. T. Gursky and J. A. Viaclovsky,
\emph{Rigidity and stability of Einstein metrics for quadratic curvature
functionals}, arXiv:1105.4648v2 (2013).
\url{https://arxiv.org/abs/1105.4648}.
The variational conventions used here are equations (2.5) and (2.9).
\bibitem{Source}
Qiushi, \emph{Einstein Gaps on Four-Manifolds: Subsequent Mathematical
Developments}, 8 October 2026,
\texttt{03\_Subsequent\_Mathematical\_Developments(1).tex}.
Harmonic-flux decomposition, equality theorem, and constrained product
variation are the supplied inputs discussed in Section 4.
\bibitem{Prior}
\emph{Global period constraints for the Q3 minimizing problem},
8 October 2026,
\texttt{Q3\_Global\_Period\_Constraints.tex}; and the earlier supplied
four-problem synthesis. These supply the fixed attainment problem, not
an existence theorem.
\end{thebibliography}
\end{document}
